\section{Certificate contents and reproduction}
The accompanying directory \nolinkurl{reproducibility/} contains the exact
input, the numerical programs, and the complete finite enclosures. Its
\nolinkurl{immutable_manifest.json} specifies SHA-256 hashes for every
unchanged numerical input, program, and certificate file. The following
dependencies distinguish finite data from mathematical hypotheses.
\begin{enumerate}
\item The exact input is \nolinkurl{head/sources/candidate_rational.json}.
Its polynomial coefficients are printed in (\ref{eq:L}) and
Appendix~\ref{app:polynomial}. Its SHA-256 hash is
\begin{quote}\small\ttfamily
207b7ed03d6bc1fd44922863e361544bb87442d13801870f0a8bc6705e8a3040.
\end{quote}
Only the exact rational coefficient strings determine the polynomials;
diagnostic fields in the file are not assumptions.
\item The moment calculation uses \nolinkurl{head/run/gauss_nodes_112.json},
\nolinkurl{head/run/head501_config.json}, and all 32 files in
\nolinkurl{head/run/blocks/}. The merged file is
\nolinkurl{head/run/head_moments_501.json}. Its full index set is
$\mathcal I$; the complete quadrature and infinite remainders are
(\ref{eq:uniform-error}) and (\ref{eq:moment-tail}). The configuration
binds the input, source, node, and block parameters. Each block is bound
to that configuration and contains exactly its prescribed four panels.
\item The two complete composed coefficient files are
\nolinkurl{head/run/forward_head_501.json} and
\nolinkurl{head/evidence/independent_head.json}. Exact rational powers
and the index proof in (\ref{eq:head-formula}) establish that no
coefficient through degree 501 is missing. Moment arithmetic and the
rational-power assembly use 384 bits; the separate interval-series
assembly uses 160 bits. The exact node containment calculation uses
rational arithmetic and does not accept approximate root locations.
\item The norm calculation uses the files
\nolinkurl{independent_integrals.json} and
\nolinkurl{independent_circle.json} in \nolinkurl{norm/evidence/}.
These contain all 266 integrals and all 512 closed panels with their
14 terms, including the domain inequalities for each estimate.
Equations (\ref{eq:S-integral})--(\ref{eq:S-tail}) justify the integral
values and their complete errors; equations
(\ref{eq:density-mass})--(\ref{eq:panel-majorant}) justify every other
finite term. The filename and the auxiliary degree-251 tail field of
the norm engine do not limit its derivative bound: its integral and
circle calculations depend only on $L$, $P$, and derivative order four.
The degree-501 tail is calculated separately from (\ref{eq:tail-final}).
\item The file \nolinkurl{rational_certificate.json} contains every
integer $k_j$ in (\ref{eq:rational-square}), the exact rational head
endpoints, and all the rational inequalities
(\ref{eq:Cmargin}) and (\ref{eq:tail-rational})--(\ref{eq:reciprocal}).
It binds the coefficient, integral, and circle files by their hashes.
The rational caps are at least both the saved panel majorants and
independently re-summed products of all their saved term enclosures.
\end{enumerate}

The arithmetic implementation uses Python, python-flint 0.8.0, and
FLINT 3.3.1. Its mathematical requirements are inclusion-preserving
real and complex ball operations, correct exact rational arithmetic,
and the certified integral algorithm described above. Floating-point
Legendre initializers, approximate optimization data, stored success
flags, and prior integral caches are not premises. SciPy is optional:
if unavailable, stored node midpoints provide initial guesses, after
which all brackets, signs, and weights are proved again.

From the directory containing \nolinkurl{verify.py}, the command
\begin{quote}\small\ttfamily
python -B -u verify.py --output /tmp/real-upper-verification
\end{quote}
checks hashes and numerical bindings, reproduces the complete moment
merge and the uniform error, verifies the nodes, reruns both coefficient
assemblies, and recomputes all smoothing integrals and circle panels
from an empty integral cache. It then checks the exact rational
conclusion. Recomputed integral balls may have different endpoints
because of adaptive subdivision or arithmetic versions; each must be a
valid enclosure and must itself give $C_4<C_*$. In particular, the
printed value of $Q$ specifies the supplied rational certificate and
does not require bitwise equality of independently recomputed
quadrature results.

The separate program \nolinkurl{reproduce.py} initializes new nodes and
quadrature errors and recomputes every moment block, with one file per
block so that interruptions can be resumed. It then performs both
assemblies, a new integral and circle calculation, and the rational
comparison. It neither installs software nor acquires computational
resources. Its command-line options and complete reproduction commands
are in \nolinkurl{README.txt} and \nolinkurl{commands.txt}. The proof of the
infinite coefficient bound is the regularity and Parseval argument in
the main text; no finite file or numerical sample substitutes for that
argument.
